%% ===========================================================================
\FloatBarrier
\section{Introduction}
\label{sec:intro}
%% ===========================================================================

Early device studies show that the electromechanical response of a member a
few micrometers thick is not the response of the same material in bulk. Bending a strip of thickness $h$ to a
curvature $\kappa_c$ produces a strain gradient of order $\kappa_c/h$, so a
coupling that is proportional to the gradient grows as $h^{-2}$ while a coupling
proportional to the strain itself does not; at the micrometer scale the two are
no longer separated by orders of magnitude, and the design of a sensor or a
transducer has to decide which of them it uses. The devices that pose the
question are small and getting smaller: vibration energy harvesters and resonant
microsensors~\cite{Roundy2003,ErturkInman2011}, mechanical metamaterials with
architected microstructure~\cite{Jiao2023}, and additively manufactured parts
whose mesostructure sets a length of its own~\cite{Sarar2024}. These early studies
establish that a device can be designed around a resonant mode of a
millimeter-scale member, measured, and reproduced by a single-length-scale
model; what they leave open is the behavior of the same member when it is
thinner than the mechanism it is built around, which is the regime in which a
gradient coupling and a dielectric memory stop being corrections. Two further
properties of the same devices matter at that scale and are absent from
classical piezoelectricity. The first is that a dielectric driven over decades
of frequency does not respond with a single time constant: its permittivity
follows a broadened dispersion, so the constitutive law is non-local in
time~\cite{ColeCole1941,HavriliakNegami1967}. The second is that the electrodes
are not a detail: whether a polarization is collected depends on which pair of
faces carries the charge, and a formulation that does not distinguish the two
configurations cannot say which effect a measurement is seeing.

Subsequent strain-gradient elasticity supplied the length that the first
property needs, and it did so in a definite order. Toupin~\cite{Toupin1962} proposed a theory of
elastic dielectrics in which the energy depends on the second displacement
gradient, Mindlin and Tiersten~\cite{MindlinTiersten1962} cast the same idea as
a theory of couple stresses, and Mindlin~\cite{Mindlin1964,Mindlin1965} extended
it to a full hierarchy of higher-order elastic constants, first for isotropic
solids and then for crystal classes. Aifantis~\cite{Aifantis1992} proposed to
reduce that hierarchy to a single parameter by keeping only the Laplacian of the
strain in the energy, and Ru and Aifantis~\cite{RuAifantis1993} extended the
reduction to a form that keeps the qualitative predictions of the full theory:
localization is regularized, singularities are removed and stiffness grows as
the member shrinks. In parallel, Eringen and
Edelen~\cite{EringenEdelen1972,Eringen1972Polar,Eringen1983} replaced the local
stress by an integral average over a neighborhood, so the length entered as the
support of a kernel rather than as a coefficient, and Lim~\cite{Lim2015}
combined the two descriptions in a higher-order nonlocal strain-gradient model
in which an integral kernel and a gradient expansion act at once. These
formulations are mathematically complete and geometrically silent: the length is
a constant of the model, and the model cannot say which feature of a
microstructure it stands for. The gap was narrowed from the other side. Reviews
by Askes and Aifantis~\cite{AskesAifantis2011} and
Polizzotto~\cite{Polizzotto2015} compare the formulations, the identification
procedures and the finite element implementations; discrete-to-continuum
arguments show how the gradient terms emerge from a lattice or a representative
volume~\cite{MetrikineAskes2002,Gitman2005}; and homogenization of
three-dimensional lattices establishes that a strain-gradient continuum
reproduces size effects that classical homogenization
misses~\cite{Molavitabrizi2023,YangLiuXia2024}. What none of these routes
delivers is the tensor itself in a form a microscopist could measure: separable
anisotropic gradient models, internal-length tensors and symmetry
classifications of three-dimensional strain-gradient elasticity provide
one~\cite{LazarPo2015,Polizzotto2018,Auffray2019}, but each of them postulates
the tensor and then fits its components to the very test it is meant to predict.
On the criterion of whether the internal length can be measured before it is
used, the averaging domain is the one route that has not been taken.

Later, flexoelectricity supplied the electromechanical side, and its history
is shorter. Kogan~\cite{Kogan1964} identified the coupling of polarization to the
strain gradient, Mindlin~\cite{Mindlin1968} formulated it as a
polarization-gradient theory, Cross~\cite{Cross2006} reviewed its magnitude in
insulating solids, Yudin and Tagantsev~\cite{YudinTagantsev2013} and Zubko and
co-workers~\cite{Zubko2013} reviewed its role in thin films, Codony and
co-workers~\cite{CodonyAriasSuryanarayana2021} computed its transverse component
from first principles for nanostructures at finite deformation, and Lam and
co-workers~\cite{Lam2003} confirmed in micron-scale bending experiments the size
dependence that the gradient predicts. On the computational side the
formulations divide by how they make the second gradient admissible. Mixed and
staggered methods carry it as an independent field and keep $C^0$
interpolation~\cite{AmanatidouAravas2002,AskesMorataAifantis2008,
PhunpengBaiz2015}; micromorphic and virtual element formulations enlarge the
kinematics~\cite{Forest2009,WriggersHudobivnik2023}; Hermite elements carry it
in the interpolation itself~\cite{BognerFoxSchmit1966,PeteraPittman1994,
Zervos2008,Papanicolopulos2009}. Flexoelectric-specific implementations have
followed both routes, with isogeometric analysis for the frequency-domain
response of flexoelectric systems~\cite{ChenYaoYvonnet2023}, computational
homogenization for their topology
optimization~\cite{ChenYvonnetYaoPark2021}, gradient electromechanical finite
elements of arbitrary order~\cite{JoshanSantapuri2024} and a strain-gradient
phase-field formulation for ferroelectrics~\cite{Wang2024PhaseField}; benchmarks
with analytical counterparts~\cite{Yang2021,Lazar2024} and isogeometric
nonlocal-strain-gradient analyses of
nanostructures~\cite{PhungVan2023EC39,PhungVan2024CS335,PhungVan2024TWS198,
Thai2023CS316} complete the landscape. The two routes are not equivalent on the
criterion that decides a coupled calculation, namely how much is spent to make
the gradient admissible: a mixed method adds a symmetric tensor field on every
element and has to stabilize it, whereas a $C^1$ element adds the second
derivatives of the deflection as nodal values and pays in continuity conditions
instead. The mixed route therefore dominates where the geometry is complex and
the mesh must follow it, and the $C^1$ route where the same operator is inverted
many times, as it is when the dielectric law carries a memory and every time
step costs a solve.

Dielectric loss is the third ingredient, and its description predates either
of the other two. A relaxing dielectric is described by a Cole--Cole or a
Havriliak--Negami dispersion~\cite{ColeCole1941,HavriliakNegami1967}, whose
time-domain counterpart is a constitutive law with a memory. Fractional calculus
gives that law one parameter more than elasticity: the spring-pot of
Caputo~\cite{Caputo1967}, used as a constitutive element by Bagley and
Torvik~\cite{BagleyTorvik1983} and by Koeller~\cite{Koeller1984}, reproduces a
power-law relaxation with the order $\alpha$ as the only addition, and it has
been carried into time-domain damping models~\cite{EnelundLesieutre1999} and
into fractional thermo-piezo-elastic wave
problems~\cite{Gupta2026Rayleigh,Gupta2024Impact,Gupta2025Piezo,
Gupta2024Vibrational}. These works establish that one extra order reproduces a
measured dispersion over decades; they leave open what the same order does to
the damping of a device whose stiffness is itself size dependent, because the
element is fitted to a dielectric measurement and then used in a mechanical
time march in which the stiffness is not the bulk one. The cost of that march is
set by where the convolution is carried rather than by the constitutive law: on
the assembled system it grows with the number of degrees of freedom and with
every step, on a small set of modal coordinates it does not, and the price is a
projection that has to be shown to be converged. A modeling decision is usually
left implicit in this literature as well: in a condensed electromechanical
problem two stiffness matrices are available, the condensed one and the
short-circuit one, and a fractional derivative multiplying one is not the same
model as the same derivative multiplying the other.

Taken together, three gaps remain, one for each line. First, the averaging window that gives the
internal length its geometric meaning has not been placed inside a coupled
electromechanical element: the coupled formulations reviewed above use a scalar
length or an assumed tensor, so the relation between the microstructure and the
length that the element sees is not stated. Second, a fractional dielectric
memory has not been formulated in a strain-gradient piezoelectric element with
an explicit statement of which matrix carries the derivative, so the effect of
that choice on a damping ratio is unknown. Third, the coupling factor of a
strain-gradient electromechanical strip has been reported for one electrode
arrangement at a time, without a measure that can be compared across thicknesses
and without separating a material property from a size effect.

Table~\ref{tab:comparison} places the present work against representative
entries from the three lines. Its rows are formulations, not implementations,
and the dashes mark ingredients that the cited paper does not carry; the last
two rows are the one-scale gradient theory of~\cite{fem1} and the present paper.
Read down the window column, the table shows that the geometric length is the
ingredient no coupled formulation carries, and read across the memory column
that the choice of memory matrix is nowhere stated.

%% The comparison table of the introduction.  One row per formulation cited in
%% the text; a dash means the ingredient is not part of that formulation.
\begin{table}[!htbp]
\centering
\small
\caption{Ingredients carried by representative strain-gradient and
flexoelectric formulations. A dash means the ingredient is not part of the
cited formulation; ``splines'' marks $C^1$ isogeometric interpolation. Only
the three last columns distinguish the present work from the one-scale theory
it extends.}
\label{tab:comparison}
\begin{tabular}{@{}lcccccc@{}}
\toprule
formulation & nonlocal or & piezo- & flexo- & frac- & $C^1$ & aniso- \\
 & gradient & electric & electric & tional & & tropic \\
 & & & & memory & & length \\
\midrule
Mindlin, second gradient~\cite{Mindlin1965}          & gradient & -- & -- & -- & -- & -- \\
Eringen, nonlocal elasticity~\cite{Eringen1972Polar} & nonlocal & -- & -- & -- & -- & -- \\
Aifantis~\cite{Aifantis1992}                         & gradient & -- & -- & -- & -- & -- \\
Polizzotto~\cite{Polizzotto2018}                     & gradient & -- & -- & -- & -- & \checkmark \\
Amanatidou and Aravas~\cite{AmanatidouAravas2002}    & gradient & -- & -- & -- & mixed & -- \\
Papanicolopulos and Zervos~\cite{Papanicolopulos2009} & gradient & -- & -- & -- & \checkmark & -- \\
Chen, Yao and Yvonnet~\cite{ChenYaoYvonnet2023}      & gradient & -- & \checkmark & -- & splines & -- \\
Joshan and Santapuri~\cite{JoshanSantapuri2024}      & gradient & -- & \checkmark & -- & -- & -- \\
one-scale theory~\cite{fem1}                         & gradient & -- & -- & -- & \checkmark & \checkmark \\
\textbf{present work} & gradient & \checkmark & \checkmark & \checkmark & \checkmark & \checkmark \\
\bottomrule
\end{tabular}
\end{table}


The novelty lies in three items, one per gap, carried by a single
discretization.

\begin{enumerate}
\item \textbf{A rectangular nonlocal window inside a coupled
  electromechanical element.} The strain-gradient energy is generated by the
  in-plane second-moment tensor $M$ of the averaging window, for which cuboid,
  ellipse and ellipsoid windows give $M_{ii}=a_i^{2}/3$, $a_i^{2}/4$ and
  $a_i^{2}/5$ in closed form. Because the energy is linear in $M$, two windows
  with the same second moment produce identical solutions: the shape of the
  window is not a parameter once $M$ is fixed, whereas its anisotropy and its
  orientation are, and they move the extensional coupling by
  $15$--$19\,\%$.
\item \textbf{Fractional dielectric memory carried on modal coordinates.} The
  dielectric relaxation is a fractional Kelvin--Voigt law of Caputo order
  $\alpha$, discretized by the L1 scheme and convoluted on the reduced modal
  coordinates, which makes the choice of memory matrix an explicit modeling
  decision; at $\alpha=1$ the short-circuit memory damps $0.450$ times as much
  as the condensed one.
\item \textbf{Mode-free coupling factors for two electrode configurations of
  one strip, and the calibration that goes with them.} The coupling factor is
  evaluated from a quasi-static expression that involves no eigenvector, so the
  same measure can be compared across geometries and thicknesses; with it the
  extensional coupling is a material property ($k^{2}=0.5581$, constant to
  $0.09\,\%$ between $10$ and $1\,\mu$m) and the flexoelectric bending coupling
  is a size effect that grows by a factor $24$ over the same range. Each of the
  four model parameters is read from a different observable, and the ranges
  over which the two-dimensional reduction and the small-coupling expansions
  hold are stated.
\end{enumerate}

Section~\ref{sec:formulation} states the formulation and derives the
constitutive relations and the governing equations,
Sec.~\ref{sec:limits} its particular cases, Sec.~\ref{sec:bc} the boundary and
the electrode conditions, and Sec.~\ref{sec:fem} the element, the condensation
and the time stepping. Section~\ref{sec:verification} establishes that the
element solves those equations, Sec.~\ref{sec:results} answers the design
questions, Sec.~\ref{sec:calibration} identifies the parameters and states the
limitations, and Sec.~\ref{sec:conclusions} collects the conclusions.

\FloatBarrier
